\section{Introdução ao Capítulo}
\label{sec:revbiblio_intro}

\paragraph{}Este capítulo apresenta a revisão bibliográfica que fundamenta o desenvolvimento do \textit{MinervaMesh}, com foco em problemas térmicos e de mecânica dos fluidos resolvidos por Métodos Numéricos, em especial o Método das Diferenças Finitas (MDF) e o Método dos Elementos Finitos (MEF). A revisão está organizada em quatro blocos: fundamentos de transferência de calor, bases dos métodos numéricos, marcos do MEF/CFD e trabalhos correlatos.

\section{Problemas Térmicos e Equações Governantes}

\subsection{Condução de Calor}
\paragraph{}A condução é descrita pela Lei de Fourier e pela equação da difusão térmica. Em regime estacionário, a formulação conduz às equações de Laplace e Poisson; em regime transiente, inclui-se o termo temporal associado ao armazenamento de energia \cite{Incropera2006,Ozisik1993,Cengel2015,Bejan1992}. Esses modelos são a base de problemas clássicos de engenharia, como dissipação em componentes eletrônicos e análise térmica de discos de freio.

\subsection{Convecção e Acoplamento com Escoamento}
\paragraph{}Quando há interação sólido-fluido, a convecção influencia a taxa de remoção de calor e exige o acoplamento entre equações de energia e de escoamento. A formulação por volumes/elementos finitos para problemas convectivo-difusivos é amplamente discutida na literatura de CFD \cite{Maliska2004,Versteeg2007,Fox2010,White2006,Bird2002}.

\subsection{Radiação Térmica}
\paragraph{}Em aplicações de alta temperatura, os termos radiativos podem se tornar relevantes no balanço térmico global \cite{Incropera2006,Cengel2015}. Embora o escopo principal deste trabalho esteja em condução e convecção, a revisão considera radiação como extensão natural para trabalhos futuros.

\section{Métodos Numéricos para EDPs}

\subsection{Método das Diferenças Finitas (MDF)}
\paragraph{}O MDF aproxima derivadas por expansões de Taylor, convertendo EDPs em sistemas algébricos. Sua implementação é direta em malhas estruturadas e, por isso, ele é frequentemente adotado em validações iniciais e em discretização temporal de problemas transientes \cite{Versteeg2007,Patankar1980}.

\paragraph{}Esquemas explícitos e implícitos (Euler e Crank-Nicolson) apresentam diferentes compromissos entre custo computacional, estabilidade e precisão. Em problemas dominados por difusão, esquemas implícitos são preferidos por robustez numérica \cite{Versteeg2007,Maliska2004}.

\subsection{Método dos Elementos Finitos (MEF)}
\paragraph{}O MEF oferece maior flexibilidade geométrica por meio de malhas não estruturadas e formulação variacional. Essa característica é decisiva em domínios complexos, comuns em aplicações reais de engenharia térmica e de fluidos \cite{Fish2007,Lewis2004,Reddy2005,Zienkiewicz2013,Bathe1996,Logan2007,Huebner2001}.

\paragraph{}A formulação de Galerkin transforma a EDP na forma fraca e permite a montagem de matrizes globais de rigidez, massa e convecção. Para elementos lineares triangulares, a implementação é eficiente e adequada a plataformas computacionais didáticas \cite{Reddy2005,Zienkiewicz2013,Fish2007}.

\section{Bases Históricas e Consolidação do MEF/CFD}

\paragraph{}A literatura mostra uma evolução contínua desde os fundamentos variacionais até formulações estabilizadas para escoamentos com convecção dominante. Trabalhos de Galerkin, Hrennikoff e Courant consolidaram a base matemática para aproximações por partes \cite{Galerkin1915,Hrennikoff1941,Courant1943}. Em seguida, contribuições de Turner e Clough estruturaram o MEF moderno aplicado à engenharia \cite{Turner1956,Clough1960}.

\paragraph{}No contexto de transporte convectivo e Navier-Stokes, formulações do tipo Petrov-Galerkin, como o SUPG, reduziram oscilações numéricas e ampliaram a faixa de estabilidade dos métodos \cite{Brooks1982,Donea1984,Christie1976,Lohner1984}. Benchmarks clássicos de cavidade e degrau regressivo permanecem referências para validação de códigos \cite{Ghia1982,Marchi2009,Thomas1981,Armaly1983,Pironneau1982}.

\paragraph{}Também se destaca a relevância de ferramentas de geração de malha na qualidade da solução, com destaque para o Gmsh em fluxos de trabalho acadêmicos e de pesquisa \cite{Geuzaine2009}.

\section{Trabalhos Correlatos}

\paragraph{}Em trabalhos de graduação recentes, observa-se forte convergência para aplicações de condução de calor e CFD com implementação numérica em Python. Na área térmica, destacam-se estudos sobre discos de freio, dissipadores e componentes eletrônicos heterogêneos \cite{Aguiar2020,Alves2021,Ferreira2022,Pinheiro2022,Capitanio2023,Miranda2023,Teixeira2023,Oliveira2025}. Esses trabalhos reforçam a importância de malhas adequadas, tratamento consistente de contornos e validação por casos de referência.

\paragraph{}Na linha de validação metodológica, a comparação entre solução analítica, MDF e MEF foi explorada por Araujo \cite{Araujo2022}, fornecendo evidências de convergência e comportamento de erro em diferentes discretizações. Em aplicações de fluidos e transferência de calor conjugada, também há contribuições relevantes com foco em escoamento para arrefecimento de eletrônicos e análise de campos de velocidade/vorticidade \cite{Florentino2022,Santos2022,SilvaJessica2020,Amaral2020,GomesMatheus2021}.

\paragraph{}No contexto local mais recente, trabalhos como os de Gomes \cite{GomesYuri2025} e Oliveira \cite{Oliveira2025} reforçam a continuidade da linha de pesquisa em transferência de calor com viés variacional/numérico e MEF tridimensional.

\section{Síntese da Revisão e Lacuna Atacada}

\paragraph{}A revisão indica que: (i) há base teórica consolidada para problemas térmicos e de escoamento; (ii) MDF e MEF são métodos complementares, com o MEF mais adequado para geometrias complexas; e (iii) existe demanda por ferramentas didáticas acessíveis para implementação e validação desses métodos.

\paragraph{}Nesse cenário, o \textit{MinervaMesh} se posiciona ao integrar formulações clássicas de engenharia numérica em uma plataforma web, com foco em acessibilidade, reprodutibilidade e uso educacional, apoiando-se em bibliotecas científicas amplamente adotadas \cite{Harris2020,PyScript2023}.
